\documentclass[tikz]{standalone}

\usepackage{tikz}
\usepackage{tkz-euclide}
\usepackage{fourier-otf}
\usepackage{fontspec}

\usetikzlibrary{calc}
\usetikzlibrary{3d}

\tikzset{
  every picture/.append style={scale=2, every node/.style={scale=2}}
}



\begin{document}
  \begin{tikzpicture}[scale=0.375]
    % Définition des limites
    \def\xmin{-4}
    \def\xmax{10}
    \def\ymin{-4}
    \def\ymax{10}
    
    % Axes
    \draw[->] (\xmin,0) -- (\xmax,0);
    \draw[->] (0,\ymin) -- (0,\ymax);
    
    % Graduations sur les axes
    \foreach \x in {\xmin,...,\xmax} {
        \ifnum\x\neq0
            \draw (\x,0) -- (\x,0.1);
        \fi
    }
    \foreach \y in {\ymin,...,\ymax} {
        \ifnum\y\neq0
            \draw (0,\y) -- (0.1,\y);
        \fi
    }
    
    % Origine
    \node[below left] at (0,0) {$O$};
    
    % Vecteurs unitaires
    \draw[thick, ->] (0,0) -- (1,0);
    \node[below] at (0.5,0) {\small $\vec{\imath}$};
    \draw[thick, ->] (0,0) -- (0,1);
    \node[left] at (0,0.5) {\small $\vec{\jmath}$};
    
    % Asymptote oblique y = x
    \draw[blue] (-11,-11) -- (9,9);
    
    % Étiquette de l'asymptote
    \node[right, rotate=45] at (6,6) {$y = x$};
    
    % Fonction f(x) = x + ln(x) pour x ≥ 1
    \draw[red, thick, domain=1:\xmax, samples=200] plot (\x, {\x + ln(\x)});
    
\end{tikzpicture}
\end{document}
